% GENERATED FILE. Edit canonical lessons, then run npm run generate:books.
% canonical-source-hash: fnv1a64:5665435bf3ae026e
% canonical-lessons: KA-C407-samsodhane, KA-C407-prayoga, KA-C407-siddhanta, KA-C407-vislesisu, KA-C407-samikse

\chapter{Doing Research}
\label{ch:ka-doing-research}
\hlchaptermodality{407}

\begin{chapteropening}
\textbf{By the end of this chapter:} \emph{I can say research, an experiment, a theory, to analyse and a survey.}

Complete the last lesson of chapter 407: I can say research, an experiment, a theory, to analyse and a survey.
\end{chapteropening}

\section[\texorpdfstring{\kn{ಸಂಶೋಧನೆ} (saṃśōdhane)}{ಸಂಶೋಧನೆ (saṃśōdhane)}]{\kn{ಸಂಶೋಧನೆ} (saṃśōdhane) — research}

\label{lesson:KA-C407-samsodhane}



\begin{quote}
Before the new one: say the Kannada for a professor, then the Kannada for to instruct.
\end{quote}

\subsection*{You'll want to know: \kn{ಸಂಶೋಧನೆ}}
\textbf{\kn{ಸಂಶೋಧನೆ}} — \emph{saṃśōdhane} — \textquotedblleft{}research\textquotedblright{}.

\begin{quote}
\kn{ಅವಳು} \kn{ಕಾಡಿನ} \kn{ಹಕ್ಕಿಗಳ} \kn{ಬಗ್ಗೆ} \kn{ಸಂಶೋಧನೆ} \kn{ಮಾಡುತ್ತಾಳೆ}.
\end{quote}

\begin{quote}
\emph{avaḷu kāḍina hakkigaḷa bagge saṃśōdhane māḍuttāḷe.}
\end{quote}

\begin{quote}
She does research on forest birds.
\end{quote}

\kn{ಸಂಶೋಧಕ}, a researcher.

\begin{cousinweb}
\textbf{\kn{ಸಂಶೋಧನೆ}} (\emph{saṃśōdhane}) is borrowed from Sanskrit \emph{saṃśodhana}, \textquotedblleft{}completely purifying, correcting\textquotedblright{} (Monier-Williams).
\end{cousinweb}

\begin{grammarlens}[title={What you've built}]
One more word for giving reasons and opinions.
\end{grammarlens}

\subsection*{Your turn}
\begin{itemize}
\raggedright
  \item \emph{Say it:} \emph{saṃśōdhane}
  \item \emph{Say it:} \emph{saṃśōdhane}, once more
  \item \emph{Say it:} \emph{avaḷu kāḍina hakkigaḷa bagge saṃśōdhane māḍuttāḷe.} — She does research on forest birds.
  \item \emph{Recall:} say the Kannada for a professor, then the Kannada for to instruct, then say \emph{saṃśōdhane} again
\end{itemize}

\begin{tcolorbox}[breakable,colback=teal!4,colframe=teal!35!black,title={Before you move on}]
What does \kn{ಸಂಶೋಧನೆ} mean? (\textquotedblleft{}Research\textquotedblright{}.) Say it once more.
\end{tcolorbox}
\section[\texorpdfstring{\kn{ಪ್ರಯೋಗ} (prayōga)}{ಪ್ರಯೋಗ (prayōga)}]{\kn{ಪ್ರಯೋಗ} (prayōga) — an experiment}

\label{lesson:KA-C407-prayoga}



\begin{quote}
Before the new one: say the Kannada for to instruct, then the Kannada for research.
\end{quote}

\subsection*{You'll want to know: \kn{ಪ್ರಯೋಗ}}
\textbf{\kn{ಪ್ರಯೋಗ}} — \emph{prayōga} — \textquotedblleft{}an experiment\textquotedblright{}.

\begin{quote}
\kn{ಮಕ್ಕಳು} \kn{ಶಾಲೆಯಲ್ಲಿ} \kn{ಸಣ್ಣ} \kn{ಪ್ರಯೋಗ} \kn{ಮಾಡಿದರು}.
\end{quote}

\begin{quote}
\emph{makkaḷu śāleyalli saṇṇa prayōga māḍidaru.}
\end{quote}

\begin{quote}
The children did a small experiment at school.
\end{quote}

\kn{ಪ್ರಯೋಗಾಲಯ}, a laboratory.

\begin{cousinweb}
\textbf{\kn{ಪ್ರಯೋಗ}} (\emph{prayōga}) is borrowed from Sanskrit \emph{prayoga}, \textquotedblleft{}joining, use, application\textquotedblright{} (Monier-Williams).
\end{cousinweb}

\begin{grammarlens}[title={What you've built}]
One more word for giving reasons and opinions.
\end{grammarlens}

\subsection*{Your turn}
\begin{itemize}
\raggedright
  \item \emph{Say it:} \emph{prayōga}
  \item \emph{Say it:} \emph{prayōga}, once more
  \item \emph{Say it:} \emph{makkaḷu śāleyalli saṇṇa prayōga māḍidaru.} — The children did a small experiment at school.
  \item \emph{Recall:} say the Kannada for to instruct, then the Kannada for research, then say \emph{prayōga} again
\end{itemize}

\begin{tcolorbox}[breakable,colback=teal!4,colframe=teal!35!black,title={Before you move on}]
What does \kn{ಪ್ರಯೋಗ} mean? (\textquotedblleft{}An experiment\textquotedblright{}.) Say it once more.
\end{tcolorbox}
\section[\texorpdfstring{\kn{ಸಿದ್ಧಾಂತ} (siddhānta)}{ಸಿದ್ಧಾಂತ (siddhānta)}]{\kn{ಸಿದ್ಧಾಂತ} (siddhānta) — a theory; a doctrine}

\label{lesson:KA-C407-siddhanta}



\begin{quote}
Before the new one: say the Kannada for research, then the Kannada for an experiment.
\end{quote}

\subsection*{You'll want to know: \kn{ಸಿದ್ಧಾಂತ}}
\textbf{\kn{ಸಿದ್ಧಾಂತ}} — \emph{siddhānta} — \textquotedblleft{}a theory; a doctrine\textquotedblright{}.

\begin{quote}
\kn{ಈ} \kn{ಸಿದ್ಧಾಂತ} \kn{ಇನ್ನೂ} \kn{ಸಾಬೀತಾಗಿಲ್ಲ}.
\end{quote}

\begin{quote}
\emph{ī siddhānta innū sābītāgilla.}
\end{quote}

\begin{quote}
This theory has not been proved yet.
\end{quote}

\kn{ಸೈದ್ಧಾಂತಿಕ}, theoretical.

\begin{cousinweb}
\textbf{\kn{ಸಿದ್ಧಾಂತ}} (\emph{siddhānta}) is borrowed from Sanskrit \emph{siddhānta}, \textquotedblleft{}an established end, a conclusion\textquotedblright{} (CDIAL 13404).
\end{cousinweb}

\begin{grammarlens}[title={What you've built}]
One more word for giving reasons and opinions.
\end{grammarlens}

\subsection*{Your turn}
\begin{itemize}
\raggedright
  \item \emph{Say it:} \emph{siddhānta}
  \item \emph{Say it:} \emph{siddhānta}, once more
  \item \emph{Say it:} \emph{ī siddhānta innū sābītāgilla.} — This theory has not been proved yet.
  \item \emph{Recall:} say the Kannada for research, then the Kannada for an experiment, then say \emph{siddhānta} again
\end{itemize}

\begin{tcolorbox}[breakable,colback=teal!4,colframe=teal!35!black,title={Before you move on}]
What does \kn{ಸಿದ್ಧಾಂತ} mean? (\textquotedblleft{}A theory; a doctrine\textquotedblright{}.) Say it once more.
\end{tcolorbox}
\section[\texorpdfstring{\kn{ವಿಶ್ಲೇಷಿಸು} (viślēṣisu)}{ವಿಶ್ಲೇಷಿಸು (viślēṣisu)}]{\kn{ವಿಶ್ಲೇಷಿಸು} (viślēṣisu) — to analyse}

\label{lesson:KA-C407-vislesisu}



\begin{quote}
Before the new one: say the Kannada for an experiment, then the Kannada for a theory.
\end{quote}

\subsection*{You'll want to know: \kn{ವಿಶ್ಲೇಷಿಸು}}
\textbf{\kn{ವಿಶ್ಲೇಷಿಸು}} — \emph{viślēṣisu} — \textquotedblleft{}to analyse\textquotedblright{}.

\begin{quote}
\kn{ನಾನು} \kn{ಪರೀಕ್ಷೆಯ} \kn{ಫಲಿತಾಂಶಗಳನ್ನು} \kn{ವಿಶ್ಲೇಷಿಸಿದೆ}.
\end{quote}

\begin{quote}
\emph{nānu parīkṣeya phalitāṃśagaḷannu viślēṣiside.}
\end{quote}

\begin{quote}
I analysed the exam results.
\end{quote}

Past (I …): \kn{ವಿಶ್ಲೇಷಿಸಿದೆ}, \emph{viślēṣiside}.

\kn{ವಿಶ್ಲೇಷಣೆ}, an analysis.

\begin{cousinweb}
\textbf{\kn{ವಿಶ್ಲೇಷಿಸು}} (\emph{viślēṣisu}) is built on Sanskrit \emph{viśleṣa}, \textquotedblleft{}loosening, separation\textquotedblright{} (Monier-Williams), with the Kannada verb ending -\emph{isu}: to take apart.
\end{cousinweb}

\begin{grammarlens}[title={What you've built}]
One more word for giving reasons and opinions.
\end{grammarlens}

\subsection*{Your turn}
\begin{itemize}
\raggedright
  \item \emph{Say it:} \emph{viślēṣisu}
  \item \emph{Say it:} \emph{viślēṣisu}, once more
  \item \emph{Say it:} \emph{nānu parīkṣeya phalitāṃśagaḷannu viślēṣiside.} — I analysed the exam results.
  \item \emph{Recall:} say the Kannada for an experiment, then the Kannada for a theory, then say \emph{viślēṣisu} again
\end{itemize}

\begin{tcolorbox}[breakable,colback=teal!4,colframe=teal!35!black,title={Before you move on}]
What does \kn{ವಿಶ್ಲೇಷಿಸು} mean? (\textquotedblleft{}To analyse\textquotedblright{}.) Say it once more.
\end{tcolorbox}
\section[\texorpdfstring{\kn{ಸಮೀಕ್ಷೆ} (samīkṣe)}{ಸಮೀಕ್ಷೆ (samīkṣe)}]{\kn{ಸಮೀಕ್ಷೆ} (samīkṣe) — a survey}

\label{lesson:KA-C407-samikse}



\begin{quote}
Before the new one: say the Kannada for a theory, then the Kannada for to analyse.
\end{quote}

\subsection*{You'll want to know: \kn{ಸಮೀಕ್ಷೆ}}
\textbf{\kn{ಸಮೀಕ್ಷೆ}} — \emph{samīkṣe} — \textquotedblleft{}a survey\textquotedblright{}.

\begin{quote}
\kn{ಸಮೀಕ್ಷೆಯಲ್ಲಿ} \kn{ನೂರು} \kn{ಕುಟುಂಬಗಳು} \kn{ಭಾಗವಹಿಸಿದವು}.
\end{quote}

\begin{quote}
\emph{samīkṣeyalli nūru kuṭumbagaḷu bhāgavahisidavu.}
\end{quote}

\begin{quote}
A hundred families took part in the survey.
\end{quote}

\kn{ಸಮೀಕ್ಷೆ} \kn{ನಡೆಸು}, to conduct a survey.

\begin{cousinweb}
\textbf{\kn{ಸಮೀಕ್ಷೆ}} (\emph{samīkṣe}) is borrowed from Sanskrit \emph{samīkṣā}, \textquotedblleft{}thorough or close inspection\textquotedblright{} (Monier-Williams).
\end{cousinweb}

\begin{grammarlens}[title={What you've built}]
One more word for giving reasons and opinions.
\end{grammarlens}

\subsection*{Your turn}
\begin{itemize}
\raggedright
  \item \emph{Say it:} \emph{samīkṣe}
  \item \emph{Say it:} \emph{samīkṣe}, once more
  \item \emph{Say it:} \emph{samīkṣeyalli nūru kuṭumbagaḷu bhāgavahisidavu.} — A hundred families took part in the survey.
  \item \emph{Recall:} say the Kannada for a theory, then the Kannada for to analyse, then say \emph{samīkṣe} again
\end{itemize}

\begin{tcolorbox}[breakable,colback=teal!4,colframe=teal!35!black,title={Before you move on}]
What does \kn{ಸಮೀಕ್ಷೆ} mean? (\textquotedblleft{}A survey\textquotedblright{}.) Say it once more.
\end{tcolorbox}
